% generated by report/sheet_figs.py from sheets/captions.json -- do not edit
\providecommand{\sheetcaption}[1]{\ifcsname sheetcap@#1\endcsname\csname sheetcap@#1\endcsname\else\texttt{#1}\fi}
\providecommand{\sheetfig}[2][\linewidth]{\begin{figure}[htbp]\centering\picsheet{#2}{#1}\caption{\sheetcaption{#2}}\label{fig:sheet:#2}\end{figure}}
\expandafter\def\csname sheetcap@s00_cover_kpis\endcsname{Workbook sheet \emph{0 Cover \& Index}, cells \texttt{A1:I19}: Cover sheet: author block and the two KPI tile rows, every number computed by an Excel formula or exported from the program.}
\expandafter\def\csname sheetcap@s02_cube_indices\endcsname{Workbook sheet \emph{2 Primitive Vertex Data}, cells \texttt{J4:O17}: Cube index buffer: 12 triangles = 36 indices; column O is an Excel formula testing the winding with n0 . ((p1-p0) x (p2-p0)) (green = CCW = front face).}
\expandafter\def\csname sheetcap@s02_sphere_cylinder\endcsname{Workbook sheet \emph{2 Primitive Vertex Data}, cells \texttt{A42:Q55}: Sphere and cylinder vertex data (first 12 vertices each); the check columns |p| = 0.5 and sqrt(x\textasciicircum{}2+y\textasciicircum{}2) = 0.5 confirm the parametrisation.}
\expandafter\def\csname sheetcap@s03_aagun\endcsname{Workbook sheet \emph{3 Object Construction}, cells \texttt{A4:Q3454}: Construction rows of the aa gatling gun (17 rows of 83 shown; row labels are the real sheet rows). Every part is one primitive with S, R (degrees) and T relative to its parent; shading model in the last column. AA Gatling gun: pedestal, yaw / pitch mount, rotating barrel cluster.}
\expandafter\def\csname sheetcap@s03_bunker\endcsname{Workbook sheet \emph{3 Object Construction}, cells \texttt{A4:Q1331}: Construction rows of the command bunker (21 rows of 417 shown; row labels are the real sheet rows). Every part is one primitive with S, R (degrees) and T relative to its parent; shading model in the last column. Command bunker (underground, 400 parts): walls, ramp, blast doors and rooms as cubes and cylinders.}
\expandafter\def\csname sheetcap@s03_flag\endcsname{Workbook sheet \emph{3 Object Construction}, cells \texttt{A4:Q5242}: Construction rows of the flag (16 rows of 61 shown; row labels are the real sheet rows). Every part is one primitive with S, R (degrees) and T relative to its parent; shading model in the last column. Flag: pole and a chain of cloth strips each yawed by a travelling wave.}
\expandafter\def\csname sheetcap@s03_gate\endcsname{Workbook sheet \emph{3 Object Construction}, cells \texttt{A4:Q5159}: Construction rows of the boom barrier (12 rows of 22 shown; row labels are the real sheet rows). Every part is one primitive with S, R (degrees) and T relative to its parent; shading model in the last column. Gate boom barrier: post, counterweight and pivoting boom with stripes.}
\expandafter\def\csname sheetcap@s03_helicopter\endcsname{Workbook sheet \emph{3 Object Construction}, cells \texttt{A4:Q3371}: Construction rows of the scout helicopter (19 rows of 110 shown; row labels are the real sheet rows). Every part is one primitive with S, R (degrees) and T relative to its parent; shading model in the last column. Scout helicopter: fuselage, main rotor and tail rotor groups.}
\expandafter\def\csname sheetcap@s03_mlrs\endcsname{Workbook sheet \emph{3 Object Construction}, cells \texttt{A4:Q3845}: Construction rows of the mlrs rocket truck (21 rows of 158 shown; row labels are the real sheet rows). Every part is one primitive with S, R (degrees) and T relative to its parent; shading model in the last column. MLRS rocket truck: cab, elevating pod frame and rocket tubes.}
\expandafter\def\csname sheetcap@s03_primitive_mix\endcsname{Workbook sheet \emph{3 Object Construction}, cells \texttt{AA3:AI25}: Primitive mix per final object: how many cylinders, spheres, cones, cubes and frusta (and how many group nodes) build each object; Excel COUNTIFS over the 4 549-row construction table.}
\expandafter\def\csname sheetcap@s03_radar\endcsname{Workbook sheet \emph{3 Object Construction}, cells \texttt{A4:Q3566}: Construction rows of the air-defence radar (16 rows of 112 shown; row labels are the real sheet rows). Every part is one primitive with S, R (degrees) and T relative to its parent; shading model in the last column. Air-defence radar: tower and rotating antenna of tiles.}
\expandafter\def\csname sheetcap@s03_rover\endcsname{Workbook sheet \emph{3 Object Construction}, cells \texttt{A4:Q2837}: Construction rows of the target rover (13 rows of 28 shown; row labels are the real sheet rows). Every part is one primitive with S, R (degrees) and T relative to its parent; shading model in the last column. Target rover: chassis, wheels, thermal engine block (heat column drives the thermal camera).}
\expandafter\def\csname sheetcap@s03_sam\endcsname{Workbook sheet \emph{3 Object Construction}, cells \texttt{A4:Q3687}: Construction rows of the sam launcher (20 rows of 121 shown; row labels are the real sheet rows). Every part is one primitive with S, R (degrees) and T relative to its parent; shading model in the last column. SAM launcher: chassis, elevating launcher frame and missile tubes.}
\expandafter\def\csname sheetcap@s03_tank_material\endcsname{Workbook sheet \emph{3 Object Construction}, cells \texttt{A4:X5000}: Tank rows, material and thermal columns: shading, procedural surface, roughness, metalness, heat (thermal-camera value), emissive, animation text.}
\expandafter\def\csname sheetcap@s03_uav\endcsname{Workbook sheet \emph{3 Object Construction}, cells \texttt{A4:Q1427}: Construction rows of the autonomous uav (22 rows of 96 shown; row labels are the real sheet rows). Every part is one primitive with S, R (degrees) and T relative to its parent; shading model in the last column. Autonomous UAV: fuselage, four motor mounts with rotor groups, gimbal yaw / pitch chain, searchlight.}
\expandafter\def\csname sheetcap@s04_inputs_blocks\endcsname{Workbook sheet \emph{4 Transform Matrices}, cells \texttt{AY4:BL22}: Worked example inputs (S, R in degrees, T looked up by component ID with INDEX / MATCH) and the translation, yaw, pitch blocks, plus Ry x Rx.}
\expandafter\def\csname sheetcap@s04_rows_error\endcsname{Workbook sheet \emph{4 Transform Matrices}, cells \texttt{A3:AW40}: Per-component rows of sheet 4 for the UAV gimbal chain: S, R, T inputs, the 3x3 rotation entries and the max-error column (Excel M versus the code's matrix).}
\expandafter\def\csname sheetcap@s04_rows_m\endcsname{Workbook sheet \emph{4 Transform Matrices}, cells \texttt{A3:AW40}: The 12 M entries computed by Excel for the gimbal chain next to the code's matrix, and the max |difference| column.}
\expandafter\def\csname sheetcap@s05_tank_chain\endcsname{Workbook sheet \emph{5 Hierarchy}, cells \texttt{A3:AI1850}: Hierarchy sheet, tank chain: tank root \textgreater{} hull \textgreater{} turret yaw \textgreater{} gun pitch \textgreater{} mantlet / barrel recoil. Depth increases by one per joint.}
\expandafter\def\csname sheetcap@s05_world_matrix\endcsname{Workbook sheet \emph{5 Hierarchy}, cells \texttt{A4:AG1850}: World matrices of the UAV gimbal chain computed by Excel: W = W\_parent x M (each W cell is a sum of products referencing the parent's row and sheet 4), next to the code's values.}
\expandafter\def\csname sheetcap@s06_charts\endcsname{Workbook sheet \emph{6 Geometry Statistics}, cells \texttt{A13:S34}: Memory charts: bytes per mesh indexed versus non-indexed (drawn from the sheet's chart definitions).}
\expandafter\def\csname sheetcap@s06_geometry_stats\endcsname{Workbook sheet \emph{6 Geometry Statistics}, cells \texttt{A4:P11}: Memory of every GPU mesh: indexed versus non-indexed, per-mesh saving and vertex reuse, and the formula vs exporter check.}
\expandafter\def\csname sheetcap@s07_attenuation\endcsname{Workbook sheet \emph{7 Lighting}, cells \texttt{A14:K32}: Distance attenuation 1/(1 + k1 d + k2 d\textasciicircum{}2) for the beacon and the searchlight, and the spotlight cone factor versus angle; all Excel formulas.}
\expandafter\def\csname sheetcap@s07_charts\endcsname{Workbook sheet \emph{7 Lighting}, cells \texttt{M14:Z50}: Attenuation and cone curves drawn from the sheet's scatter charts.}
\expandafter\def\csname sheetcap@s08_comparison\endcsname{Workbook sheet \emph{8 Shading}, cells \texttt{A12:J31}: Same bunker, same camera: flat, Gouraud and Phong override side by side (renders from the program).}
\expandafter\def\csname sheetcap@s08_shading\endcsname{Workbook sheet \emph{8 Shading}, cells \texttt{A4:H13}: Shading models used in the scene: parts using each (COUNTIF), which normal is used, where lighting is evaluated, cost and visual effect.}
\expandafter\def\csname sheetcap@s09_features\endcsname{Workbook sheet \emph{9 Render Pipeline}, cells \texttt{A5:H17}: Per-feature frame cost: configuration, frame time, baseline, feature cost and share (differences computed by Excel).}
\expandafter\def\csname sheetcap@s09_pipeline_passes\endcsname{Workbook sheet \emph{9 Render Pipeline}, cells \texttt{A20:D30}: GPU pass list of one frame (measured): time per pass, share of the total and cumulative time, Excel formulas.}
\expandafter\def\csname sheetcap@s10_testing\endcsname{Workbook sheet \emph{10 Testing}, cells \texttt{A1:G30}: Test sheet: KPI row (tests, passed, failed, pass rate - Excel formulas over the status column) and the first 25 tests with expected and observed results.}
\expandafter\def\csname sheetcap@s14_aagun\endcsname{Workbook sheet \emph{14 Controllable Objects}, cells \texttt{A173:L650}: Controllable AA-1 GATLING: the moving parts and the animation rule of each articulated node, as exported from the code.}
\expandafter\def\csname sheetcap@s14_derived\endcsname{Workbook sheet \emph{14 Controllable Objects}, cells \texttt{A156:D169}: Derived figures computed by Excel from the specification table (e.g. turret 360 degrees = 360 / slew rate).}
\expandafter\def\csname sheetcap@s14_doors\endcsname{Workbook sheet \emph{14 Controllable Objects}, cells \texttt{A173:L650}: Controllable BLAST DOORS: the moving parts and the animation rule of each articulated node, as exported from the code.}
\expandafter\def\csname sheetcap@s14_gate\endcsname{Workbook sheet \emph{14 Controllable Objects}, cells \texttt{A173:L650}: Controllable GATE BARRIER: the moving parts and the animation rule of each articulated node, as exported from the code.}
\expandafter\def\csname sheetcap@s14_helicopter\endcsname{Workbook sheet \emph{14 Controllable Objects}, cells \texttt{A173:L650}: Controllable MH-6 HELICOPTER: the moving parts and the animation rule of each articulated node, as exported from the code.}
\expandafter\def\csname sheetcap@s14_mlrs\endcsname{Workbook sheet \emph{14 Controllable Objects}, cells \texttt{A173:L650}: Controllable MLRS-1 ROCKET TRUCK: the moving parts and the animation rule of each articulated node, as exported from the code.}
\expandafter\def\csname sheetcap@s14_radar\endcsname{Workbook sheet \emph{14 Controllable Objects}, cells \texttt{A173:L650}: Controllable AIR-DEFENCE RADAR: the moving parts and the animation rule of each articulated node, as exported from the code.}
\expandafter\def\csname sheetcap@s14_sam\endcsname{Workbook sheet \emph{14 Controllable Objects}, cells \texttt{A173:L650}: Controllable SAM-1 LAUNCHER: the moving parts and the animation rule of each articulated node, as exported from the code.}
\expandafter\def\csname sheetcap@s14_specs_tank\endcsname{Workbook sheet \emph{14 Controllable Objects}, cells \texttt{A24:F60}: Specification rows of the tank read from the code constants, with the Excel-derived figures in sheet 14 below them.}
\expandafter\def\csname sheetcap@s14_tank\endcsname{Workbook sheet \emph{14 Controllable Objects}, cells \texttt{A173:L650}: Controllable MBT-1 TANK: the moving parts and the animation rule of each articulated node, as exported from the code.}
\expandafter\def\csname sheetcap@s14_uav\endcsname{Workbook sheet \emph{14 Controllable Objects}, cells \texttt{A173:L650}: Controllable UAV-1 DRONE: the moving parts and the animation rule of each articulated node, as exported from the code.}
\expandafter\def\csname sheetcap@s16_effects\endcsname{Workbook sheet \emph{16 Combat \& Effects}, cells \texttt{A43:G54}: Particle effects measured by running the real Combat code: particle count, peak, lifetime and dynamic light for each effect.}
\expandafter\def\csname sheetcap@s16_weapons\endcsname{Workbook sheet \emph{16 Combat \& Effects}, cells \texttt{A35:I40}: Weapons: speeds, blast power, damage radius (= power x kDamageRadiusPerPower), flight time, range, reload and guidance; formulas over the constants table.}
\expandafter\def\csname sheetcap@s18_params_footprint\endcsname{Workbook sheet \emph{18 Computer Vision}, cells \texttt{A40:E59}: Derived camera figures (pixels per degree, horizontal FOV, gate radius) and the ground footprint of one pixel versus UAV altitude, all Excel formulas.}
\expandafter\def\csname sheetcap@s18_tests\endcsname{Workbook sheet \emph{18 Computer Vision}, cells \texttt{A63:E73}: Vision tests T130-T137 with their PASS status looked up from the testing sheet (INDEX / MATCH).}
\expandafter\def\csname sheetcap@s19_armour\endcsname{Workbook sheet \emph{19 Squad \& Mission}, cells \texttt{A102:D117}: Armour and damage falloff: damage per round, rounds to kill with one or two UAVs and the reload of one UAV, Excel formulas.}
\expandafter\def\csname sheetcap@s19_combined\endcsname{Workbook sheet \emph{19 Squad \& Mission}, cells \texttt{A65:G71}: Combined strike: launch delay of the nearer UAV so both rounds arrive together (arrival spread = 0).}
\expandafter\def\csname sheetcap@s19_roles\endcsname{Workbook sheet \emph{19 Squad \& Mission}, cells \texttt{A4:C16}: Squad roles and strike phases.}
\expandafter\def\csname sheetcap@s19_waves\endcsname{Workbook sheet \emph{19 Squad \& Mission}, cells \texttt{A74:D100}: Mission constants: waves, base integrity and score parameters from the code.}
\expandafter\def\csname sheetcap@s20_hud_tests\endcsname{Workbook sheet \emph{20 Audio \& HUD}, cells \texttt{A56:E71}: Audio tests with PASS status looked up from sheet 10.}
\expandafter\def\csname sheetcap@s20_mixer\endcsname{Workbook sheet \emph{20 Audio \& HUD}, cells \texttt{A24:C35}: Mixer: sample rate, voices, panning, distance gain, air absorption, Doppler.}
\expandafter\def\csname sheetcap@s20_sounds\endcsname{Workbook sheet \emph{20 Audio \& HUD}, cells \texttt{A4:D21}: Synthesised sounds: trigger, technique and key parameters of all 15 sounds (no sample files).}
\expandafter\def\csname sheetcap@sheet_construction_tank\endcsname{Workbook sheet \emph{3 Object Construction}, cells \texttt{A4:Q3261}: Construction rows of the main battle tank (14 rows of 284 shown; row labels are the real sheet rows). Every part is one primitive with S, R (degrees) and T relative to its parent; shading model in the last column. Main battle tank: hull, tracks, turret yaw, gun pitch and recoil group.}
\expandafter\def\csname sheetcap@sheet_controllables\endcsname{Workbook sheet \emph{14 Controllable Objects}, cells \texttt{A3:L21}: All controllable objects: root, part counts per primitive, group nodes and animated nodes (Excel SUMIFS over the construction sheet).}
\expandafter\def\csname sheetcap@sheet_cv\endcsname{Workbook sheet \emph{18 Computer Vision}, cells \texttt{A4:E16}: Computer-vision pipeline: the 11 stages from the adaptive threshold to the Kalman track life cycle, with the code location of each.}
\expandafter\def\csname sheetcap@sheet_cvnoise\endcsname{Workbook sheet \emph{23 CV Noise \& Seeds}, cells \texttt{A1:H60}: CV Noise \& Seeds: sensor noise configuration, P\_d / false-alarm sweep against NETD, derived stream seeds and the v6 tests with the status looked up from the testing sheet.}
\expandafter\def\csname sheetcap@sheet_daycycle\endcsname{Workbook sheet \emph{21 Day Cycle \& Lighting}, cells \texttt{A1:W40}: Day Cycle \& Lighting: the exported 24 h grid (sun elevation, key intensity, ambient, fog, exposure, solar, beams, waves, cloud-shadow mean) with the Excel recomputation of the sun elevation from latitude, declination and hour, the difference and the OK check.}
\expandafter\def\csname sheetcap@sheet_hierarchy\endcsname{Workbook sheet \emph{5 Hierarchy}, cells \texttt{A3:AI1850}: Hierarchy sheet, UAV gimbal chain: UAV root \textgreater{} gimbal yaw \textgreater{} gimbal pitch \textgreater{} lamp / lens / camera; world position and the Excel-vs-code check, with the animation text of each joint.}
\expandafter\def\csname sheetcap@sheet_lighting\endcsname{Workbook sheet \emph{7 Lighting}, cells \texttt{A4:R8}: The three scene lights plus the ambient term: type, position, direction, colour, intensity, attenuation coefficients (k linear, k quadratic), spot cone angles, role.}
\expandafter\def\csname sheetcap@sheet_memory\endcsname{Workbook sheet \emph{2 Primitive Vertex Data}, cells \texttt{J20:Q39}: Technique 1 (separate triangles) versus Technique 2 (unique vertices + indices) for one face and for the whole cube; every figure is an Excel formula and the last row checks it against the exporter's byte count.}
\expandafter\def\csname sheetcap@sheet_mission\endcsname{Workbook sheet \emph{19 Squad \& Mission}, cells \texttt{A40:F62}: Time-on-target model of the rail rounds: ramp distance, flight time t(d) and its closed-form check against roundFlightTime(d); the last row confirms formula = code.}
\expandafter\def\csname sheetcap@sheet_mmult\endcsname{Workbook sheet \emph{4 Transform Matrices}, cells \texttt{BD11:BQ29}: Worked example for one component (CAMP-55): rotation blocks Ry, Rx, Rz, then R = Ry x Rx x Rz, T x R and finally M = T x R x S, every cell an Excel sum of products; the last block is |M - M\_code|.}
\expandafter\def\csname sheetcap@sheet_soldiers\endcsname{Workbook sheet \emph{22 Soldiers \& Routes}, cells \texttt{A1:Q40}: Soldiers \& Routes: roster with day / night duty, primitive counts per soldier (Excel totals), route lengths, joint tree and gait constants.}
\expandafter\def\csname sheetcap@sheet_vertex_cube\endcsname{Workbook sheet \emph{2 Primitive Vertex Data}, cells \texttt{A4:H29}: Cube vertex buffer: 24 unique vertices (4 per face) as interleaved x, y, z, nx, ny, nz; the face is derived from the normal by an Excel formula. (The lab primitives carry no texture coordinates.)}
\providecommand{\autocaption}[2]{\ifcsname figcap@#1@#2\endcsname\csname figcap@#1@#2\endcsname\else\texttt{#2}\fi}
\providecommand{\autofig}[3][\linewidth]{\begin{figure}[htbp]\centering\pic{#2/#3}{#1}\caption{\autocaption{#2}{#3}}\label{fig:#2:#3}\end{figure}}
\expandafter\def\csname figcap@diagrams@A1_cube_vertices\endcsname{The unit cube generator: 24 vertices (4 per face, each with the face normal) and 36 indices (2 counter-clockwise triangles per face, diagonal 0-2). Indexed storage needs 720 B instead of 864 B. {\footnotesize(Source: \texttt{src/primitives/Cube.cpp: createCube()})}}
\expandafter\def\csname figcap@diagrams@A2_sphere_rings\endcsname{The sphere generator: 614 vertices on 9 rings x 36 slices per hemisphere, 1224 triangles; the normal of every vertex equals its normalised position, the southern half is the mirrored northern half. {\footnotesize(Source: \texttt{src/primitives/Sphere.cpp: createSphere(36, 18)})}}
\expandafter\def\csname figcap@diagrams@A3_cylinder_cone\endcsname{Cylinder, cone and frustum come from one generator: two interleaved rings (bottom z = -0.5, top z = +0.5), side triangles, and a triangle fan per cap; a cone is simply top radius 0, the beam frustum has no caps. {\footnotesize(Source: \texttt{src/primitives/Cylinder.cpp: createCylinder(); Renderer.cpp lines 136-141 (which primitives exist)})}}
\expandafter\def\csname figcap@diagrams@A4_terrain_disc\endcsname{The terrain is a 36 m disc of 49921 vertices (130 rings x 384 slices) built like a cylinder cap and displaced by a height function; the left panel shows the ring structure on a coarse copy, the right panel the real triangles with the entrance-ramp hole. {\footnotesize(Source: \texttt{src/objects/Terrain.cpp: createTerrainDisc(), terrainHeight()})}}
\expandafter\def\csname figcap@diagrams@A5_trs_worked_example\endcsname{Worked example for one real part, the tank's left turret cheek: the unit cube is scaled, rotated about Y and translated into its parent (the turret); the printed matrices are the scene's own, and T*R*S reproduces the node's local matrix. {\footnotesize(Source: \texttt{src/scene/Transform.h: Transform::matrix(); src/math/Math3D.cpp: rotateEulerYXZ(); Tank.cpp (node 'Turret cheek L')})}}
\expandafter\def\csname figcap@diagrams@A6_hierarchy_tank\endcsname{Scene-graph tree of the main battle tank (collapsed: identical siblings are merged): hull -\textgreater{} track groups / turret yaw -\textgreater{} gun pitch -\textgreater{} barrel recoil -\textgreater{} muzzle; each group multiplies its local matrix onto its children. {\footnotesize(Source: \texttt{src/objects/Tank.cpp (rig), src/scene/SceneNode.h; counts from the live Scene})}}
\expandafter\def\csname figcap@diagrams@A6_hierarchy_uav\endcsname{Scene-graph trees of the UAV (gimbal yaw -\textgreater{} gimbal pitch carries lamp and thermal camera; rotors under their motor mounts; rounds on the launcher rails) and of the intruder rover. {\footnotesize(Source: \texttt{src/objects/Uav.cpp, src/objects/Rover.cpp; counts from the live Scene})}}
\expandafter\def\csname figcap@diagrams@B1_barycentric\endcsname{Rasteriser interpolation: every pixel gets barycentric weights of its triangle; colours (Gouraud) or normals (Phong) are mixed with them, shown for a red / green / blue triangle. The right plot shows why the two differ: lighting is not linear in the normal. {\footnotesize(Source: \texttt{OpenGL rasteriser; ShaderSources.cpp Gouraud / Phong varyings (TS\_TERM\_VARYINGS, line \textasciitilde{}652)})}}
\expandafter\def\csname figcap@diagrams@B2_shading_models_coarse\endcsname{The same lit sphere mesh (sphere small) shaded three ways and compared with the exact analytic sphere; RMSE flat 0.143, Gouraud 0.118, Phong 0.112. Lighting uses the shader's own brdf() ported to Python. {\footnotesize(Source: \texttt{ShaderSources.cpp: brdf() 113-139, computeLight() 143-196, compose() 561-566, kFlatGeom \textasciitilde{}676; ShadersPostFx.cpp acesFitted() 387})}}
\expandafter\def\csname figcap@diagrams@B2_shading_models_real\endcsname{The same lit sphere mesh (sphere) shaded three ways and compared with the exact analytic sphere; RMSE flat 0.050, Gouraud 0.039, Phong 0.037. Lighting uses the shader's own brdf() ported to Python. {\footnotesize(Source: \texttt{ShaderSources.cpp: brdf() 113-139, computeLight() 143-196, compose() 561-566, kFlatGeom \textasciitilde{}676; ShadersPostFx.cpp acesFitted() 387})}}
\expandafter\def\csname figcap@diagrams@B3_lambert\endcsname{Lambert's cosine law as used for the diffuse term: intensity falls with cos of the angle between normal and light and is clamped at the terminator; the sphere at the right is lit by this term only. {\footnotesize(Source: \texttt{ShaderSources.cpp brdf() lines 113-139 (diff = radiance * nl)})}}
\expandafter\def\csname figcap@diagrams@C1_specular_lobes\endcsname{Surface shininess as a lobe: Blinn-Phong power n (v1) vs the GGX distribution (PBR); the right plot maps each v1 shininess to its derived roughness (roughness = 1 - log2(n)/8): GGX has a similar core width but a longer tail, which gives the soft halo around highlights. {\footnotesize(Source: \texttt{ShaderSources.cpp brdf() 119-137; Material.h pbrRoughness() line 61})}}
\expandafter\def\csname figcap@diagrams@C2_fresnel_geometry\endcsname{The G and F factors of the Cook-Torrance BRDF: Smith-Schlick shadowing, Schlick Fresnel for a dielectric (F0 = 0.04) and for the UAV's metal parts, and the resulting specular response when the view sweeps through the mirror direction for three roughness values. {\footnotesize(Source: \texttt{ShaderSources.cpp brdf() lines 127-137, computeLight() line 148 (F0 = mix(0.04, albedo, metalness))})}}
\expandafter\def\csname figcap@diagrams@C3_specular_grid\endcsname{Roughness x metalness grid of lit spheres: low roughness gives a small bright highlight, high roughness a wide dim one; with metalness the diffuse term vanishes (1 - metalness) and F0 takes the albedo, so the highlight is tinted. {\footnotesize(Source: \texttt{ShaderSources.cpp brdf() 113-139, computeLight() 143-150, compose() 561-566})}}
\expandafter\def\csname figcap@diagrams@C4_practical_attenuation\endcsname{Practical lights fall off like an inverse square but are multiplied by a smooth window so they reach exactly zero at their range; spot lights add a smoothstep cone. The histogram lists the ranges of the lights actually built into the scene. {\footnotesize(Source: \texttt{ShaderSources.cpp practicalAttenuation() lines 90-100; ranges from Scene::practicalLights()})}}
\expandafter\def\csname figcap@diagrams@D1_ssr_raymarch\endcsname{How a puddle mirrors the scene: the reflected ray is stepped in world space, every sample is projected to the screen and its depth compared with the depth buffer; the first crossing (within a thickness) is refined by six bisections and the scene colour at that pixel is added, faded near the screen edge and for long marches. {\footnotesize(Source: \texttt{ShadersPostFx.cpp kCombineFrag reflection() lines 44-81 (constants 0.04, 1.09, 0.03, thickness 0.15 + 0.04 t + 0.01 d, 6 bisections)})}}
\expandafter\def\csname figcap@diagrams@D2_fresnel_reflectivity\endcsname{Reflection strength: the Fresnel term makes grazing views mirror-like while the composite keeps at least 35 \%; the lit pass decides per surface how much reflection it asks for (wet top faces, polished metal, puddles). {\footnotesize(Source: \texttt{ShadersPostFx.cpp lines 75-77, 91; ShaderSources.cpp compose() 576-582, terrainShade() 541-546})}}
\expandafter\def\csname figcap@diagrams@E1_planner_grids\endcsname{The four fields the infiltrating rover plans on: the free-cell occupancy grid, the exposure to the two UAVs, the per-rover noise field of two seeds and the resulting step cost (all dumped from the live HideController of the scene). {\footnotesize(Source: \texttt{src/ai/HideController.cpp: setStaticOccluders(), computeVisibility(), noiseAt(), floodFillToObjective()})}}
\expandafter\def\csname figcap@diagrams@E2_flood_dijkstra\endcsname{Weighted MicroMouse flood field (Dijkstra from the objective) for two seeds with the same start and objective: the cost-to-go heat map, the steepest-descent cell route and the smoothed guide line; different noise fields make different routes. {\footnotesize(Source: \texttt{src/ai/HideController.cpp: floodFillToObjective() lines 213-247, descend() 144-172, replan() 262-282})}}
\expandafter\def\csname figcap@diagrams@E3_flood_bfs\endcsname{The v4 hiding behaviour: cells hidden from the UAV become goals, a breadth-first flood fill gives every free cell its step distance to cover, and the rover walks down the field (MicroMouse). Compared with the Dijkstra field, all steps cost the same. {\footnotesize(Source: \texttt{src/ai/HideController.cpp: computeVisibility(), floodFill() lines 120-142, descend()})}}
\expandafter\def\csname figcap@diagrams@E4_path_chaikin\endcsname{Chaikin corner cutting turns the staircase of grid cells into the smooth guide line the line follower tracks; three passes at 1/4 - 3/4 points, shown step by step on the real route. {\footnotesize(Source: \texttt{src/ai/HideController.cpp: chaikin() lines 28-41, replan() lines 262-282})}}
\expandafter\def\csname figcap@diagrams@E5_ir_sensors\endcsname{The line follower: five IR reflectance sensors on a bar 0.31 m ahead of the robot read whether they are over the 10 cm guide line; the weighted mean of the active offsets is the error that drives a PD controller and the differential-drive wheel speeds. Drawn from a real frame of a simulated drive. {\footnotesize(Source: \texttt{src/ai/HideController.cpp: followLine() lines 308-366; constants kSensorForward, kSensorSpacing, kLineHalfWidth, kKp, kKd})}}
\expandafter\def\csname figcap@diagrams@E6_rover_trace\endcsname{A recorded drive of the intruder rover: its path over the occupancy grid and, for a 22 s window, the sensor bits, line error, PD steering command and left / right wheel speeds (error and PD command are reconstructed from the recorded sensor bits with the controller's formulas). {\footnotesize(Source: \texttt{src/ai/HideController.cpp followLine(); data from Scene::update() with figdata drive})}}
\expandafter\def\csname figcap@diagrams@F10_run_timeline\endcsname{Timeline of a simulated 90 s defence: the two UAVs' roles, the strike phase, the tracks' position uncertainty by life-cycle state, the game events (spotting, lock, launch, combined strike, explosion, kill) and the number of live rovers. {\footnotesize(Source: \texttt{src/scene/Scene.cpp update(); src/ai/Squad.cpp (data: figdata dumpSquad, 90 s, 30 Hz, vision from tests/test\_squad.cpp syntheticBlobs)})}}
\expandafter\def\csname figcap@diagrams@F11_mission_waves\endcsname{Mission waves: wave n sends min(2 + n, 6) intruders (3, 4, 5, then 6), with the scoring and breach rules; the bars are produced by running the real Mission class. {\footnotesize(Source: \texttt{src/scene/Mission.cpp: startWave() lines 24-35, update() 37-75 (data: figdata dumpMission)})}}
\expandafter\def\csname figcap@diagrams@F1_threshold_histogram\endcsname{Adaptive threshold of the image stage: the histogram of a heat frame with the fixed floor 0.80, Otsu's threshold (maximum of the between-class variance) and the threshold actually used, their maximum. A hot frame raises Otsu above the floor; a cold one leaves the floor in charge. {\footnotesize(Source: \texttt{src/vision/CvPipeline.cpp: otsuThreshold() lines 9-32, detectTargets() lines 114-122 (synthetic 320x240 frame through the real functions)})}}
\expandafter\def\csname figcap@diagrams@F2_cv_stages\endcsname{The image stage on a synthetic 320x240 heat frame, run through the real functions: threshold, despeckle, labelling, features and the target filter; the table lists why each component is accepted or rejected. {\footnotesize(Source: \texttt{src/vision/CvPipeline.cpp: detectTargets() lines 106-174, despeckle() 34-52, labelComponents() 66-114})}}
\expandafter\def\csname figcap@diagrams@F3_union_find\endcsname{Two-pass connected-component labelling with union-find on a small 12 x 7 mask: pass 1 gives provisional labels and records that two of them touch (red pixels); pass 2 resolves every label to its root. The picture is checked against the real labelComponents() output. {\footnotesize(Source: \texttt{src/vision/CvPipeline.cpp: labelComponents() lines 66-114, findRoot() 56-64})}}
\expandafter\def\csname figcap@diagrams@F4_blob_moments\endcsname{Blob features from image moments: centroid, bounding box and the principal axes of the pixel distribution (eigenvectors of the second central moments) give the orientation and the elongation used to reject thin streaks; values were recomputed and checked against the real detectTargets(). {\footnotesize(Source: \texttt{src/vision/CvPipeline.cpp: detectTargets() lines 126-158})}}
\expandafter\def\csname figcap@diagrams@F5_kalman_filter\endcsname{The track filter on a rover driving straight, turning and driving straight: noisy ground measurements, the Kalman estimate with 2-sigma ellipses, the position error against the filter's own uncertainty, and the velocity estimate converging from the unknown initial value. Run through the real MultiTracker. {\footnotesize(Source: \texttt{src/vision/MultiTracker.cpp: predict() lines 21-46, kalmanUpdate() 60-82 (data: figdata dumpKalman)})}}
\expandafter\def\csname figcap@diagrams@F6_gating_hungarian\endcsname{Data association for one UAV frame: each track predicts where it should be, the chi-square gate (d\textasciicircum{}2 \textless{} 9.21) removes impossible pairs, and the Hungarian algorithm picks the pairing with the smallest total Mahalanobis cost (note m0 and m1 both fit track 3 but only the globally cheaper pairing wins). {\footnotesize(Source: \texttt{src/vision/MultiTracker.cpp: mahalanobis2() 48-56, hungarian() 84-137, update() 175-190 (data: figdata dumpAssociation)})}}
\expandafter\def\csname figcap@diagrams@F7_triangulation\endcsname{Two-view triangulation: each UAV's pixel becomes a ray, the fused fix is the midpoint of the shortest segment between the rays (residual = segment length); Monte Carlo over pointing noise shows it beats a single ground-plane ray and that a wing UAV at about 90 degrees gives the best geometry. {\footnotesize(Source: \texttt{src/vision/MultiTracker.cpp: triangulate() lines 139-153; Squad.cpp measurement sigma 0.12 + 0.015 r (data: figdata dumpTriangulation)})}}
\expandafter\def\csname figcap@diagrams@F8_squad_state_machine\endcsname{The squad's state machines: the UAV roles (patrol, track, lead / wing, search and back) and the strike phases (tracking, planned, inbound, reacquire), with the lock times and distances that trigger the transitions. {\footnotesize(Source: \texttt{src/ai/Squad.cpp: update() 220-327, planStrike() 176-218; Squad.h constants kSquadLeadLock etc.})}}
\expandafter\def\csname figcap@diagrams@F9_time_on_target\endcsname{Time-on-target: the squad predicts each round's flight time from the distance with a two-phase model (accelerating launch, then 60 m/s, detonating 0.5 m short), launches the farther UAV first and the nearer one later so the blasts coincide; the right plot shows the model against the simulated missile and the distances of this strike. {\footnotesize(Source: \texttt{src/ai/Squad.cpp: roundFlightTime() lines 43-50, planStrike() 193-212; src/scene/Combat.h kMissile* (data: figdata dumpSquad / dumpFlight)})}}
\expandafter\def\csname figcap@diagrams@G1_render_pipeline\endcsname{Overview of one rendered frame: scene graph to world matrices, shadow maps and pre-pass, ambient occlusion and fog, the HDR object pass with Flat / Gouraud / Phong lighting, then SSR, temporal anti-aliasing, exposure, bloom / depth of field and the final filmic composite (or the thermal sensor model). {\footnotesize(Source: \texttt{src/rendering/RendererFrame.cpp (frame order), ShadersPostFx.cpp header comment lines 1-4})}}
\expandafter\def\csname figcap@diagrams@G2_cv_pipeline\endcsname{From heat pixels to a shared track picture and a coordinated strike: image stage (threshold, despeckle, labelling, features, target filter), ground ray, gating and assignment, Kalman update with optional two-view triangulation, track life cycle, and the squad logic. {\footnotesize(Source: \texttt{src/vision/CvPipeline.h header comment, MultiTracker.h header comment, Squad.h header comment})}}
\expandafter\def\csname figcap@diagrams@G3_game_loop\endcsname{One simulation step: input, vision blobs, squad decisions, controllable objects, mission waves, intruder planning and driving, combat and effects, then world matrices and rendering; the rendered thermal images become the next frame's vision input. {\footnotesize(Source: \texttt{src/scene/Scene.cpp: Scene::update() lines 353-445; src/app/App.cpp update() lines 262-285})}}
\expandafter\def\csname figcap@diagrams@H1_daycycle_curves\endcsname{Day cycle from the exported 24 h grid (97 rows): sun elevation with the phase thresholds and the four phase events, and the lighting quantities that follow it (key intensity, daylight weight, practical lamps, exposure key, mean cloud shadow); bottom: the key-light hue (sun reddening by air mass, moon blue at night). {\footnotesize(Source: \texttt{build/scene\_export.json v6.day\_cycle.grid (DayCycle::applyTo, src/scene/DayCycle.cpp)})}}
\expandafter\def\csname figcap@diagrams@H2_kasten_young\endcsname{Left: the Kasten--Young air mass against the plane-parallel 1/sin(e), which diverges at the horizon. Right: the per-channel sun transmittance relative to the zenith path and the normalised disc colour (strip): white at noon, orange-red near the horizon. {\footnotesize(Source: \texttt{src/scene/DayCycle.cpp airMass / sunTransmittance (ported), kSunOpticalDepth from DayCycle.h})}}
\expandafter\def\csname figcap@diagrams@H3_catmull_rom\endcsname{Left: Catmull--Rom curves through the same unevenly spaced control points; the uniform parameterisation overshoots next to the short edges, the centripetal one stays tight without cusps. Right: curvature along the centripetal curve and the walking speed it allows (soldier walk speed 1.3 m/s). {\footnotesize(Source: \texttt{src/math/Spline.cpp evalSegment / curvatureAtDistance and src/scene/Garrison.cpp stepRoute (ported); control points illustrative})}}
\expandafter\def\csname figcap@diagrams@H4_gait_cycle\endcsname{The walk cycle as the program computes it: joint angles over one cycle (left stance shaded), the ankle path relative to the hip (flat world-locked stance, cycloid swing with lift) and stick figures rebuilt by forward kinematics from the same angles. {\footnotesize(Source: \texttt{build/scene\_export.json v6.garrison.gait (gaitPose, footInPelvis, legIK in src/scene/Gait.cpp)})}}
\expandafter\def\csname figcap@diagrams@H5_two_link_ik\endcsname{Two-link IK by the law of cosines: the knee angle from the triangle (hip, knee, ankle), the hip angle from the direction to the target plus the triangle's angle at the hip. The same solver (with a pole vector instead of the planar knee rule) poses the arms on the rifle. {\footnotesize(Source: \texttt{src/scene/Gait.cpp legIK / legFK (ported); thigh and shin lengths from v6.garrison.gait.walk})}}
\expandafter\def\csname figcap@diagrams@H6_spawn_ring\endcsname{Random intruder routes: a spawn bearing drawn uniformly on the ring outside the gate-road gap, up to three via points in a sector around the spawn bearing (driven outermost first, so the chain closes in on the camp), then the objective. Right: the rejection sampler gives a flat bearing histogram with an empty gap at north. {\footnotesize(Source: \texttt{src/scene/Mission.cpp randomRoute (ported, Python RNG); constants from v6.random\_routes and BaseLayout.h})}}
\expandafter\def\csname figcap@diagrams@H7_pd_netd\endcsname{Real-CV statistics against the temporal noise level (NETD), full sensor model (PSF, fixed pattern after NUC, stuck pixels): detection probability of a rover (full size and half size), person recognition, false targets and persons misread as targets per frame; crosses: the same at the default NETD without NUC. Bottom: the synthetic test crop at each NETD with the classified blobs (target = cyan, person = green). {\footnotesize(Source: \texttt{build/figdata\_out/cv\_noise.json (tools/figdata\_main.cpp dumpCvNoise; SensorModel + cv::detectAll)})}}
\expandafter\def\csname figcap@diagrams@attenuation\endcsname{Practical lights fall off like an inverse square but are multiplied by a smooth window so they reach exactly zero at their range; spot lights add a smoothstep cone. The histogram lists the ranges of the lights actually built into the scene. {\footnotesize(Source: \texttt{ShaderSources.cpp practicalAttenuation() lines 90-100; ranges from Scene::practicalLights()})}}
\expandafter\def\csname figcap@diagrams@flood_bfs\endcsname{The v4 hiding behaviour: cells hidden from the UAV become goals, a breadth-first flood fill gives every free cell its step distance to cover, and the rover walks down the field (MicroMouse). Compared with the Dijkstra field, all steps cost the same. {\footnotesize(Source: \texttt{src/ai/HideController.cpp: computeVisibility(), floodFill() lines 120-142, descend()})}}
\expandafter\def\csname figcap@diagrams@flood_dijkstra\endcsname{Weighted (Dijkstra) flood field of the infiltrating rover with the cost-to-go values, the steepest-descent route and the smoothed guide line (seed 11). {\footnotesize(Source: \texttt{src/ai/HideController.cpp: floodFillToObjective(), descend()})}}
\expandafter\def\csname figcap@diagrams@fresnel\endcsname{The G and F factors of the Cook-Torrance BRDF: Smith-Schlick shadowing, Schlick Fresnel for a dielectric (F0 = 0.04) and for the UAV's metal parts, and the resulting specular response when the view sweeps through the mirror direction for three roughness values. {\footnotesize(Source: \texttt{ShaderSources.cpp brdf() lines 127-137, computeLight() line 148 (F0 = mix(0.04, albedo, metalness))})}}
\expandafter\def\csname figcap@diagrams@frustum_ndc\endcsname{The perspective projection: the view frustum of the main camera (fov 50 deg, near 0.05 m, far 600 m), the matrix built by perspective() with its actual numbers, and the resulting non-linear NDC depth, which is why the shaders linearise depth before using it. {\footnotesize(Source: \texttt{src/math/Math3D.cpp: perspective() lines 84-94; src/scene/Camera.h: nearZ = 0.05, farZ = 600, fovDeg = 50})}}
\expandafter\def\csname figcap@diagrams@gating_ellipse\endcsname{Data association for one UAV frame: each track predicts where it should be, the chi-square gate (d\textasciicircum{}2 \textless{} 9.21) removes impossible pairs, and the Hungarian algorithm picks the pairing with the smallest total Mahalanobis cost (note m0 and m1 both fit track 3 but only the globally cheaper pairing wins). {\footnotesize(Source: \texttt{src/vision/MultiTracker.cpp: mahalanobis2() 48-56, hungarian() 84-137, update() 175-190 (data: figdata dumpAssociation)})}}
\expandafter\def\csname figcap@diagrams@ggx_curves\endcsname{Surface shininess as a lobe: Blinn-Phong power n (v1) vs the GGX distribution (PBR); the right plot maps each v1 shininess to its derived roughness (roughness = 1 - log2(n)/8): GGX has a similar core width but a longer tail, which gives the soft halo around highlights. {\footnotesize(Source: \texttt{ShaderSources.cpp brdf() 119-137; Material.h pbrRoughness() line 61})}}
\expandafter\def\csname figcap@diagrams@ir_sensors\endcsname{The line follower: five IR reflectance sensors on a bar 0.31 m ahead of the robot read whether they are over the 10 cm guide line; the weighted mean of the active offsets is the error that drives a PD controller and the differential-drive wheel speeds. Drawn from a real frame of a simulated drive. {\footnotesize(Source: \texttt{src/ai/HideController.cpp: followLine() lines 308-366; constants kSensorForward, kSensorSpacing, kLineHalfWidth, kKp, kKd})}}
\expandafter\def\csname figcap@diagrams@kalman_demo\endcsname{The track filter on a rover driving straight, turning and driving straight: noisy ground measurements, the Kalman estimate with 2-sigma ellipses, the position error against the filter's own uncertainty, and the velocity estimate converging from the unknown initial value. Run through the real MultiTracker. {\footnotesize(Source: \texttt{src/vision/MultiTracker.cpp: predict() lines 21-46, kalmanUpdate() 60-82 (data: figdata dumpKalman)})}}
\expandafter\def\csname figcap@diagrams@path_chaikin\endcsname{Chaikin corner cutting turns the staircase of grid cells into the smooth guide line the line follower tracks; three passes at 1/4 - 3/4 points, shown step by step on the real route. {\footnotesize(Source: \texttt{src/ai/HideController.cpp: chaikin() lines 28-41, replan() lines 262-282})}}
\expandafter\def\csname figcap@diagrams@raster_edges\endcsname{How the GPU decides which pixels a triangle covers: three edge functions are evaluated at every pixel centre; a pixel is inside where all three are non-negative, and the normalised values are exactly the barycentric weights used to interpolate colours, normals and depth. {\footnotesize(Source: \texttt{OpenGL rasteriser (conceptual); weights used by the Gouraud / Phong varyings in ShaderSources.cpp})}}
\expandafter\def\csname figcap@diagrams@rover_trace\endcsname{A recorded drive of the intruder rover: its path over the occupancy grid and, for a 22 s window, the sensor bits, line error, PD steering command and left / right wheel speeds (error and PD command are reconstructed from the recorded sensor bits with the controller's formulas). {\footnotesize(Source: \texttt{src/ai/HideController.cpp followLine(); data from Scene::update() with figdata drive})}}
\expandafter\def\csname figcap@diagrams@ssr_raymarch\endcsname{How a puddle mirrors the scene: the reflected ray is stepped in world space, every sample is projected to the screen and its depth compared with the depth buffer; the first crossing (within a thickness) is refined by six bisections and the scene colour at that pixel is added, faded near the screen edge and for long marches. {\footnotesize(Source: \texttt{ShadersPostFx.cpp kCombineFrag reflection() lines 44-81 (constants 0.04, 1.09, 0.03, thickness 0.15 + 0.04 t + 0.01 d, 6 bisections)})}}
\expandafter\def\csname figcap@diagrams@triangulation\endcsname{Two-view triangulation: each UAV's pixel becomes a ray, the fused fix is the midpoint of the shortest segment between the rays (residual = segment length); Monte Carlo over pointing noise shows it beats a single ground-plane ray and that a wing UAV at about 90 degrees gives the best geometry. {\footnotesize(Source: \texttt{src/vision/MultiTracker.cpp: triangulate() lines 139-153; Squad.cpp measurement sigma 0.12 + 0.015 r (data: figdata dumpTriangulation)})}}
\expandafter\def\csname figcap@figures_lab@buffers_overview\endcsname{The render pipeline buffers of one frame: moon shadow maps, depth, normals, SSAO, fog in-scatter and transmittance, HDR colour and its false-colour log luminance, bloom, reflectivity, final. {\footnotesize(Rendered: \texttt{glGetTexImage of the renderer's internal textures after endFrame (RendererLab.cpp), mapped to 8-bit for display.})}}
\expandafter\def\csname figcap@figures_lab@bump_on_off\endcsname{Surface relief OFF (top) / ON (bottom) on six materials under raking light. {\footnotesize(Rendered: \texttt{Copy of relief\_compare.png (bumpOn flag -\textgreater{} uBumpOff).})}}
\expandafter\def\csname figcap@figures_lab@cv_despeckle\endcsname{Despeckle result, removed pixels in red. {\footnotesize(Rendered: \texttt{Copy of cv\_uN\_c\_despeckle.png.})}}
\expandafter\def\csname figcap@figures_lab@cv_f_external\endcsname{(f) The same simulation frame from an external camera: both UAVs, the confirmed track with its uncertainty, and each UAV's pixel ray to the target. {\footnotesize(Rendered: \texttt{Renderer::beginFrame/endFrame on the live scene, Hud brackets from projectMarker, tracks from the squad's MultiTracker.})}}
\expandafter\def\csname figcap@figures_lab@cv_frame\endcsname{Measurement numbers of the frame: per UAV Otsu / threshold / counts, per blob area, centroid, orientation, elongation, ground point, track id; per track state, sigma, triangulation residual and error against the rover's true position. {\footnotesize(Rendered: \texttt{Written by lab::cvFrame.})}}
\expandafter\def\csname figcap@figures_lab@cv_labels\endcsname{Connected components, a random colour per label. {\footnotesize(Rendered: \texttt{Copy of cv\_uN\_d\_components.png.})}}
\expandafter\def\csname figcap@figures_lab@cv_pipeline_uav1\endcsname{Computer-vision stage on UAV 1's real 320x240 thermal frame: raw heat (ironbow and grey), histogram with the Otsu threshold, threshold mask, despeckle, connected components, centroid / box / principal axes, and the same frame seen from outside. {\footnotesize(Rendered: \texttt{Rendered by the program (--capture --only lab\_cv): raw heat from Renderer::captureUavCamera (view HeatRaw), stages from cv::detectTargetsDebug (the same functions as detectTargets), overlays with the Hud API. Moment chosen automatically: first frame with both UAVs holding a confirmed track.})}}
\expandafter\def\csname figcap@figures_lab@cv_pipeline_uav2\endcsname{The same pipeline on UAV 2's frame at the same instant. {\footnotesize(Rendered: \texttt{As cv\_pipeline\_uav1.png.})}}
\expandafter\def\csname figcap@figures_lab@cv_strip\endcsname{Six frames, 0.15 s apart: the detection box follows the rover; white trail = centroid history. {\footnotesize(Rendered: \texttt{Frames captured while stepping the simulation; cv::detectTargets on each; track id = nearest track to the blob's ground point.})}}
\expandafter\def\csname figcap@figures_lab@cv_thermal\endcsname{Real raw 320x240 heat image of the UAV that sees the rover best, x3, ironbow. {\footnotesize(Rendered: \texttt{Copy of cv\_uN\_a\_raw\_ironbow.png.})}}
\expandafter\def\csname figcap@figures_lab@cv_threshold\endcsname{Threshold mask of that frame with the threshold value printed. {\footnotesize(Rendered: \texttt{Copy of cv\_uN\_b\_mask.png.})}}
\expandafter\def\csname figcap@figures_lab@cv_tracks\endcsname{Tracker output overlaid on the same frame seen from outside: tracks, uncertainty, UAV pixel rays. {\footnotesize(Rendered: \texttt{Copy of cv\_f\_external.png.})}}
\expandafter\def\csname figcap@figures_lab@cv_u1_a_raw_grey\endcsname{(a) Raw heat image of UAV 1 in grey (white-hot palette). {\footnotesize(Rendered: \texttt{As above, palette 1.})}}
\expandafter\def\csname figcap@figures_lab@cv_u1_a_raw_ironbow\endcsname{(a) Raw heat image of UAV 1, 320x240 scaled x3 (nearest), ironbow palette. {\footnotesize(Rendered: \texttt{Hud::heatImage of the R8 heat texture, nearest filter.})}}
\expandafter\def\csname figcap@figures_lab@cv_u1_b0_histogram\endcsname{Histogram (log counts) of the heat image with the fixed floor 204, the Otsu level and the threshold actually used. {\footnotesize(Rendered: \texttt{cv::otsuThreshold + Hud rectangles.})}}
\expandafter\def\csname figcap@figures_lab@cv_u1_b_mask\endcsname{(b) Threshold mask; the printed value is t = max(floor 204, Otsu + 1). {\footnotesize(Rendered: \texttt{Pixels \textgreater{}= CvFrameStats::threshold.})}}
\expandafter\def\csname figcap@figures_lab@cv_u1_c_despeckle\endcsname{(c) After despeckle: pixels with fewer than 2 hot 8-neighbours are removed (red). {\footnotesize(Rendered: \texttt{cv::despeckle on the mask; removed = mask and not cleaned.})}}
\expandafter\def\csname figcap@figures_lab@cv_u1_d_components\endcsname{(d) Connected components (two-pass labelling, union-find, 8-connectivity), one random colour per label. {\footnotesize(Rendered: \texttt{cv::labelComponents; colour = hash of the label.})}}
\expandafter\def\csname figcap@figures_lab@cv_u1_e_overlay\endcsname{(e) Accepted blobs: bounding box (cyan), centroid (white cross), principal axes from the second central moments (yellow = major, magenta = minor, length 2 sigma); rejected components in red. {\footnotesize(Rendered: \texttt{Moments recomputed from the labelled pixels; the target filter is CvConfig (area, hot core, elongation).})}}
\expandafter\def\csname figcap@figures_lab@cv_u2_a_raw_grey\endcsname{(a) UAV 2 raw heat image (grey). {\footnotesize(Rendered: \texttt{As cv\_u1\_a\_raw\_grey.png.})}}
\expandafter\def\csname figcap@figures_lab@cv_u2_a_raw_ironbow\endcsname{(a) UAV 2 raw heat image (x3, ironbow). {\footnotesize(Rendered: \texttt{As cv\_u1\_a\_raw\_ironbow.png.})}}
\expandafter\def\csname figcap@figures_lab@material_grid\endcsname{PBR roughness x metalness grid plus Blinn-Phong row. {\footnotesize(Rendered: \texttt{Copy of materials\_grid.png.})}}
\expandafter\def\csname figcap@figures_lab@materials_blinnphong_row\endcsname{Blinn-Phong row: shininess 4, 16, 64, 256. {\footnotesize(Rendered: \texttt{uPbr = 0.})}}
\expandafter\def\csname figcap@figures_lab@materials_grid\endcsname{Material grid: 5 x 5 spheres, roughness 0.1-0.9 (columns) x metalness 0 / 0.25 / 0.5 / 0.75 / 1 (rows), copper albedo, lit by a moon-like directional light and one point lamp; below, the Blinn-Phong row (shininess 4 / 16 / 64 / 256). {\footnotesize(Rendered: \texttt{Real PBR shader (GGX / Cook-Torrance, uPbr = 1) on a LabStage of unit spheres; the Blinn-Phong row uses RenderSettings::pbr = false. HIGH quality, no fog.})}}
\expandafter\def\csname figcap@figures_lab@materials_pbr_grid\endcsname{The 5 x 5 PBR grid alone. {\footnotesize(Rendered: \texttt{As materials\_grid.png.})}}
\expandafter\def\csname figcap@figures_lab@pipeline_buffers\endcsname{Internal render buffers of one frame. {\footnotesize(Rendered: \texttt{Copy of buffers\_overview.png (scaled to 1600 px).})}}
\expandafter\def\csname figcap@figures_lab@reflections_compare\endcsname{Wet ground in rain: SSR off, SSR on, the reflectivity buffer (surface target alpha) and the amplified difference. {\footnotesize(Rendered: \texttt{RenderSettings::ssrOn on the live scene (wetness 1, rain 1, ULTRA); buffer read with Renderer::labReadBuffer.})}}
\expandafter\def\csname figcap@figures_lab@relief_compare\endcsname{Procedural surface relief OFF (top) and ON (bottom) on corrugated metal, steel plate, HESCO, concrete, rock and asphalt under raking light. {\footnotesize(Rendered: \texttt{RenderSettings::bumpOn -\textgreater{} uniform uBumpOff in surfaceBump(); cubes on a LabStage, HIGH quality.})}}
\expandafter\def\csname figcap@figures_lab@shading_models\endcsname{UAV with flat / Gouraud / Phong override and wireframe. {\footnotesize(Rendered: \texttt{Copy of shading\_uav\_compare.png (scaled to 1600 px).})}}
\expandafter\def\csname figcap@figures_lab@shading_sphere_compare\endcsname{A coarse 10x6 sphere (top row) and the stock 36x18 sphere (bottom row): flat facets, Gouraud Mach bands and a lost / smeared highlight, Phong, and the triangles. {\footnotesize(Rendered: \texttt{Renderer::labSphereResolution swaps the sphere mesh; Blinn-Phong terms (pbr = false) so the highlight shows the interpolation difference.})}}
\expandafter\def\csname figcap@figures_lab@shading_uav_compare\endcsname{The UAV (hierarchical composite of unit primitives) with the shading override flat / Gouraud / Phong and a hidden-line wireframe of its triangles. {\footnotesize(Rendered: \texttt{RenderSettings::shading = Flat / Gouraud / Phong, isolate = Autonomous UAV; wire: view Diagnostic + EdgeMode::HiddenLine.})}}
\expandafter\def\csname figcap@figures_lab@shadow_cascades\endcsname{Moon shadow map near cascade and world cascade (depth). {\footnotesize(Rendered: \texttt{glGetTexImage of the 2-layer shadow array.})}}
\expandafter\def\csname figcap@figures_lab@ssr_on_off\endcsname{Screen-space reflections off / on, reflectivity buffer, difference. {\footnotesize(Rendered: \texttt{Copy of reflections\_compare.png (scaled to 1600 px).})}}
\expandafter\def\csname figcap@figures_lab@surface_patterns\endcsname{Procedural albedo patterns of the surface materials (relief off). {\footnotesize(Rendered: \texttt{Copy of relief\_off.png.})}}
